% GENERATED FILE. Edit canonical lessons, then run npm run generate:books.
% canonical-source-hash: fnv1a64:c891d5cf3186c4c7
% canonical-lessons: JA-C115-futa, JA-C115-mato, JA-C115-hinata, JA-C115-namako, JA-C115-ayu

\chapter{Lid, Target, Sunny spot, Sea cucumber, Sweetfish}
\label{ch:ja-lid-target-sunny-spot-sea-cucumber-sweetfish}
\hlchaptermodality{115}

\begin{chapteropening}
\textbf{By the end of this chapter:} \emph{I can say a lid, a target, a sunny spot, a sea cucumber and a sweetfish.}

Complete the last lesson of chapter 115: i can say a lid, a target, a sunny spot, a sea cucumber and a sweetfish.
\end{chapteropening}

\section[\texorpdfstring{\ja{ふた} (futa)}{ふた (futa)}]{\ja{ふた} (futa) — a lid}

\label{lesson:JA-C115-futa}



\begin{quote}
Before the new one: say the Japanese for honey, then the Japanese for a net.
\end{quote}

\subsection*{You'll want to know: \ja{ふた}}
\textbf{\ja{ふた}} — \emph{futa} — \textquotedblleft{}a lid\textquotedblright{}.

\begin{grammarlens}[title={What you've built}]
One more word for this chapter.
\end{grammarlens}

\subsection*{Your turn}
\begin{itemize}
\raggedright
  \item \emph{Say it:} \emph{futa}
  \item \emph{Say it:} \emph{futa}, once more
  \item \emph{Say it:} say \emph{futa}
  \item \emph{Recall:} say the Japanese for honey, then the Japanese for a net, then say \emph{futa} again
\end{itemize}

\begin{tcolorbox}[breakable,colback=teal!4,colframe=teal!35!black,title={Before you move on}]
What does \ja{ふた} mean? (\textquotedblleft{}A lid\textquotedblright{}.) Say it once more.
\end{tcolorbox}
\section[\texorpdfstring{\ja{まと} (mato)}{まと (mato)}]{\ja{まと} (mato) — a target}

\label{lesson:JA-C115-mato}



\begin{quote}
Before the new one: say the Japanese for a net, then the Japanese for a lid.
\end{quote}

\subsection*{You'll want to know: \ja{まと}}
\textbf{\ja{まと}} — \emph{mato} — \textquotedblleft{}a target\textquotedblright{}.

\begin{grammarlens}[title={What you've built}]
One more word for this chapter.
\end{grammarlens}

\subsection*{Your turn}
\begin{itemize}
\raggedright
  \item \emph{Say it:} \emph{mato}
  \item \emph{Say it:} \emph{mato}, once more
  \item \emph{Say it:} say \emph{mato}
  \item \emph{Recall:} say the Japanese for a net, then the Japanese for a lid, then say \emph{mato} again
\end{itemize}

\begin{tcolorbox}[breakable,colback=teal!4,colframe=teal!35!black,title={Before you move on}]
What does \ja{まと} mean? (\textquotedblleft{}A target\textquotedblright{}.) Say it once more.
\end{tcolorbox}
\section[\texorpdfstring{\ja{ひなた} (hinata)}{ひなた (hinata)}]{\ja{ひなた} (hinata) — a sunny spot}

\label{lesson:JA-C115-hinata}



\begin{quote}
Before the new one: say the Japanese for a lid, then the Japanese for a target.
\end{quote}

\subsection*{You'll want to know: \ja{ひなた}}
\textbf{\ja{ひなた}} — \emph{hinata} — \textquotedblleft{}a sunny spot\textquotedblright{}.

\begin{grammarlens}[title={What you've built}]
One more word for this chapter.
\end{grammarlens}

\subsection*{Your turn}
\begin{itemize}
\raggedright
  \item \emph{Say it:} \emph{hinata}
  \item \emph{Say it:} \emph{hinata}, once more
  \item \emph{Say it:} say \emph{hinata}
  \item \emph{Recall:} say the Japanese for a lid, then the Japanese for a target, then say \emph{hinata} again
\end{itemize}

\begin{tcolorbox}[breakable,colback=teal!4,colframe=teal!35!black,title={Before you move on}]
What does \ja{ひなた} mean? (\textquotedblleft{}A sunny spot\textquotedblright{}.) Say it once more.
\end{tcolorbox}
\section[\texorpdfstring{\ja{なまこ} (namako)}{なまこ (namako)}]{\ja{なまこ} (namako) — a sea cucumber}

\label{lesson:JA-C115-namako}



\begin{quote}
Before the new one: say the Japanese for a target, then the Japanese for a sunny spot.
\end{quote}

\subsection*{You'll want to know: \ja{なまこ}}
\textbf{\ja{なまこ}} — \emph{namako} — \textquotedblleft{}a sea cucumber\textquotedblright{}.

\begin{grammarlens}[title={What you've built}]
One more word for this chapter.
\end{grammarlens}

\subsection*{Your turn}
\begin{itemize}
\raggedright
  \item \emph{Say it:} \emph{namako}
  \item \emph{Say it:} \emph{namako}, once more
  \item \emph{Say it:} say \emph{namako}
  \item \emph{Recall:} say the Japanese for a target, then the Japanese for a sunny spot, then say \emph{namako} again
\end{itemize}

\begin{tcolorbox}[breakable,colback=teal!4,colframe=teal!35!black,title={Before you move on}]
What does \ja{なまこ} mean? (\textquotedblleft{}A sea cucumber\textquotedblright{}.) Say it once more.
\end{tcolorbox}
\section[\texorpdfstring{\ja{あゆ} (ayu)}{あゆ (ayu)}]{\ja{あゆ} (ayu) — a sweetfish}

\label{lesson:JA-C115-ayu}



\begin{quote}
Before the new one: say the Japanese for a sunny spot, then the Japanese for a sea cucumber.
\end{quote}

\subsection*{You'll want to know: \ja{あゆ}}
\textbf{\ja{あゆ}} — \emph{ayu} — \textquotedblleft{}a sweetfish\textquotedblright{}.

\begin{grammarlens}[title={What you've built}]
One more word for this chapter.
\end{grammarlens}

\subsection*{Your turn}
\begin{itemize}
\raggedright
  \item \emph{Say it:} \emph{ayu}
  \item \emph{Say it:} \emph{ayu}, once more
  \item \emph{Say it:} say \emph{ayu}
  \item \emph{Recall:} say the Japanese for a sunny spot, then the Japanese for a sea cucumber, then say \emph{ayu} again
\end{itemize}

\begin{tcolorbox}[breakable,colback=teal!4,colframe=teal!35!black,title={Before you move on}]
What does \ja{あゆ} mean? (\textquotedblleft{}A sweetfish\textquotedblright{}.) Say it once more.
\end{tcolorbox}
